\par\Needspace{6\baselineskip}\medskip
\hypertarget{native-111-article082}{}\pdfbookmark[2]{82.}{article-082}
\begin{center}82.\end{center}
Suppose finally that three radius vectors $r$, $r'$, $r''$ are given that make angles $N$, $N'$, $N''$ with an arbitrary straight line drawn from the Sun in the plane of the orbit. Retaining the other notation, we then have the equations
\[
\begin{aligned}
&\frac{p}{r}=1+e\cos(N-\Pi)\\[4pt]
\mathrm{I.}\qquad&\frac{p}{r'}=1+e\cos(N'-\Pi)\\[4pt]
&\frac{p}{r''}=1+e\cos(N''-\Pi),
\end{aligned}
\]
from which $p$, $\Pi$, and $e$ can be obtained in several ways. To calculate $p$ before the other quantities, multiply the three equations I respectively by $\sin(N''-N')$, $-\sin(N''-N)$, and $\sin(N'-N)$. Adding the products and applying Lemma I of Article 78 gives
\[
p=\frac{\sin(N''-N')-\sin(N''-N)+\sin(N'-N)}{\dfrac1r\sin(N''-N')-\dfrac1{r'}\sin(N''-N)+\dfrac1{r''}\sin(N'-N)}\mathpunct{\cdot}
\]
This expression deserves closer examination. The numerator clearly becomes
\[
\begin{aligned}
&=2\sin\tfrac12(N''-N')\cos\tfrac12(N''-N')\\
&\qquad-2\sin\tfrac12(N''-N')\cos(\tfrac12N''+\tfrac12N'-N)\\
&=4\sin\tfrac12(N''-N')\sin\tfrac12(N''-N)\sin\tfrac12(N'-N).
\end{aligned}
\]
If we further set
\[
\begin{aligned}
r'r''\sin(N''-N')&=n\\
rr''\sin(N''-N)&=n'\\
rr'\sin(N'-N)&=n'',
\end{aligned}
\]
it is clear that $\tfrac12n$, $\tfrac12n'$, and $\tfrac12n''$ are the areas of the triangles formed by the second and third radius vectors, by the first and third, and by the first and second, respectively. From this it is easy to see that in the new formula
\SourceLeaf{112}
\[
p=\frac{4\sin\tfrac12(N''-N')\sin\tfrac12(N''-N)\sin\tfrac12(N'-N)\mathbin{.}rr'r''}{n-n'+n''}
\]
the denominator is twice the area of the triangle with vertices at the endpoints of the three radius vectors, that is, at the three positions of the celestial body in space. Whenever these positions are close to one another, this area will always be very small, indeed of the third order if $N'-N$, $N''-N'$ are regarded as small quantities of the first order. It also follows that if one or more of the quantities $r$, $r'$, $r''$, $N$, $N'$, $N''$ are affected by errors, however slight, a very large error can result in the determination of $p$. Thus this method of finding the dimensions of the orbit can never yield high precision unless the three heliocentric positions are separated by considerable distances.

Once the semiparameter $p$ has been found, $e$ and $\Pi$ are determined by combining any two of equations I according to the method of Article 79.
